% =============================================================================
% Lecture 1: Economic Decision-Making & Everyday Choices
% =============================================================================
\documentclass[10pt,english,aspectratio=169]{beamer}
% =============================================================================
% Pathways: Economics & Finance Workshop — Beamer Preamble
% =============================================================================
% Usage: Each lecture .tex file should start with:
%   \documentclass[10pt,english,aspectratio=169]{beamer}
%   % =============================================================================
% Pathways: Economics & Finance Workshop — Beamer Preamble
% =============================================================================
% Usage: Each lecture .tex file should start with:
%   \documentclass[10pt,english,aspectratio=169]{beamer}
%   % =============================================================================
% Pathways: Economics & Finance Workshop — Beamer Preamble
% =============================================================================
% Usage: Each lecture .tex file should start with:
%   \documentclass[10pt,english,aspectratio=169]{beamer}
%   \input{../../codes/Preambles/header_slides}
% Compile: cd output/Slides && TEXINPUTS=../../codes/Preambles:$TEXINPUTS xelatex file.tex
% =============================================================================

% ---------- Packages ----------
% Fonts
\usepackage{lmodern}
\usepackage[utf8]{inputenc}
\usepackage[normalem]{ulem}

% Math
\usepackage{amsmath,amssymb,amsthm}

% Figures and tables
\usepackage{graphicx}
\usepackage{array,tabularx,booktabs,threeparttable}

% TikZ and pgfplots
\usepackage{pgfplots}
\pgfplotsset{compat=1.18}
\usepackage{tikz}
\usetikzlibrary{positioning,arrows.meta,calc}

% Colored boxes
\usepackage[most]{tcolorbox}

% Miscellaneous
\usepackage{transparent}
\usepackage{setspace}
\usepackage[nodayofweek]{datetime}
\usepackage{hyperref}

% ---------- Theme ----------
\usetheme{CambridgeUS}
\usecolortheme{dolphin}
\usepackage{appendixnumberbeamer}

% ---------- Itemize ----------
\setbeamercovered{invisible}
\setbeamertemplate{itemize items}[default]
\setbeamertemplate{enumerate items}[default]

\setbeamertemplate{itemize item}[triangle]
\setbeamertemplate{itemize subitem}[triangle]
\setbeamertemplate{itemize subsubitem}{-}

% ---------- Headline / Footline ----------
\setbeamertemplate{headline}{}

\setbeamertemplate{footline}{
  \leavevmode
  \hbox{
    \begin{beamercolorbox}[wd=\paperwidth,ht=2.5ex,dp=1.25ex,leftskip=1em,rightskip=1em]{footline}
      \hfill \insertframenumber
    \end{beamercolorbox}
  }
}

\setbeamertemplate{navigation symbols}{}

% ---------- Frame title (centered) ----------
\makeatletter
\setbeamertemplate{frametitle}{
  \nointerlineskip
  \begin{beamercolorbox}[wd=\paperwidth,sep=0.35cm,center]{frametitle}
    {\usebeamerfont{frametitle}\insertframetitle\par}
    \ifx\insertframesubtitle\@empty\else
      \vspace{0.15cm}
      {\usebeamerfont{framesubtitle}\usebeamercolor[fg]{framesubtitle}\insertframesubtitle\par}
    \fi
  \end{beamercolorbox}
}
\makeatother

% ---------- Margins ----------
\setbeamersize{
  text margin left=2em,
  text margin right=2em
}

% ---------- Section outline frames ----------
\AtBeginSection[]{
  \begin{frame}{Outline}
    \tableofcontents[currentsection,
                     currentsubsection,
                     hideothersubsections]
  \end{frame}
}

% =============================================================================
% Colors — Okabe-Ito colorblind-friendly palette
% =============================================================================

\definecolor{black}{HTML}{000000}
\definecolor{orange}{HTML}{E69F00}
\definecolor{skyblue}{HTML}{56B4E9}
\definecolor{green}{HTML}{009E73}
\definecolor{yellow}{HTML}{F0E442}
\definecolor{blue}{HTML}{0072B2}
\definecolor{vermillion}{HTML}{D55E00}
\definecolor{purple}{HTML}{CC79A7}

% Semantic aliases (mapped to Okabe-Ito)
\definecolor{softbg}{HTML}{F5F5F5}
\definecolor{lightgray}{gray}{0.65}

% =============================================================================
% Box Environments
% =============================================================================

% --- Key points (orange background) ---
\newtcolorbox{keybox}{
  colback=orange!12,
  colframe=orange,
  fonttitle=\bfseries,
  boxrule=0.5pt,
  arc=2pt,
  left=6pt,
  right=6pt,
  top=4pt,
  bottom=4pt
}

% --- Highlights (orange left accent) ---
\newtcolorbox{highlightbox}{
  colback=white,
  colframe=orange,
  fonttitle=\bfseries,
  boxrule=0pt,
  leftrule=3pt,
  arc=0pt,
  left=6pt,
  right=6pt,
  top=4pt,
  bottom=4pt
}

% --- Definitions (blue border, titled) ---
\newtcolorbox{definitionbox}[1][]{
  colback=blue!5,
  colframe=blue,
  fonttitle=\bfseries,
  title={#1},
  boxrule=0.5pt,
  arc=2pt,
  left=6pt,
  right=6pt,
  top=4pt,
  bottom=4pt
}

% --- Methods (sky blue background) ---
\newtcolorbox{methodbox}{
  colback=skyblue!8,
  colframe=skyblue!70,
  fonttitle=\bfseries,
  boxrule=0.5pt,
  arc=2pt,
  left=6pt,
  right=6pt,
  top=4pt,
  bottom=4pt
}

% --- Results (green border) ---
\newtcolorbox{resultbox}{
  colback=green!5,
  colframe=green,
  fonttitle=\bfseries,
  boxrule=0.5pt,
  arc=2pt,
  left=6pt,
  right=6pt,
  top=4pt,
  bottom=4pt
}

% --- Quotes (purple left rule) ---
\newtcolorbox{quotebox}{
  colback=softbg,
  colframe=purple,
  fonttitle=\itshape,
  boxrule=0pt,
  leftrule=2pt,
  arc=0pt,
  left=6pt,
  right=6pt,
  top=4pt,
  bottom=4pt
}

% =============================================================================
% Math Commands
% =============================================================================

\newcommand{\E}{\mathrm{E}}
\newcommand{\Var}{\mathrm{Var}}
\newcommand{\Cov}{\mathrm{Cov}}
\newcommand{\R}{\mathbb{R}}
\newcommand{\suminfty}[1]{\sum_{#1=0}^\infty}
\newcommand{\Lcal}{\mathcal{L}}
\newcommand{\prtl}[2]{\frac{\partial #1}{\partial #2}}
\renewcommand{\div}[2]{\frac{d #1}{d #2}}
\newcommand{\suchthat}{\text{ s.t. }}
\newcommand{\wrt}{\text{ w.r.t. }}
\newcommand{\red}[1]{\textcolor{vermillion}{#1}}
\newcommand{\blue}[1]{\textcolor{blue}{#1}}

% =============================================================================
% Economics Shortcuts
% =============================================================================

\newcommand{\OC}{\text{Opportunity Cost}}
\newcommand{\DWL}{\text{DWL}}
\newcommand{\pmark}{\textcolor{resultgreen}{\checkmark}}
\newcommand{\xmark}{\textcolor{darkred}{\texttimes}}

% Compile: cd output/Slides && TEXINPUTS=../../codes/Preambles:$TEXINPUTS xelatex file.tex
% =============================================================================

% ---------- Packages ----------
% Fonts
\usepackage{lmodern}
\usepackage[utf8]{inputenc}
\usepackage[normalem]{ulem}

% Math
\usepackage{amsmath,amssymb,amsthm}

% Figures and tables
\usepackage{graphicx}
\usepackage{array,tabularx,booktabs,threeparttable}

% TikZ and pgfplots
\usepackage{pgfplots}
\pgfplotsset{compat=1.18}
\usepackage{tikz}
\usetikzlibrary{positioning,arrows.meta,calc}

% Colored boxes
\usepackage[most]{tcolorbox}

% Miscellaneous
\usepackage{transparent}
\usepackage{setspace}
\usepackage[nodayofweek]{datetime}
\usepackage{hyperref}

% ---------- Theme ----------
\usetheme{CambridgeUS}
\usecolortheme{dolphin}
\usepackage{appendixnumberbeamer}

% ---------- Itemize ----------
\setbeamercovered{invisible}
\setbeamertemplate{itemize items}[default]
\setbeamertemplate{enumerate items}[default]

\setbeamertemplate{itemize item}[triangle]
\setbeamertemplate{itemize subitem}[triangle]
\setbeamertemplate{itemize subsubitem}{-}

% ---------- Headline / Footline ----------
\setbeamertemplate{headline}{}

\setbeamertemplate{footline}{
  \leavevmode
  \hbox{
    \begin{beamercolorbox}[wd=\paperwidth,ht=2.5ex,dp=1.25ex,leftskip=1em,rightskip=1em]{footline}
      \hfill \insertframenumber
    \end{beamercolorbox}
  }
}

\setbeamertemplate{navigation symbols}{}

% ---------- Frame title (centered) ----------
\makeatletter
\setbeamertemplate{frametitle}{
  \nointerlineskip
  \begin{beamercolorbox}[wd=\paperwidth,sep=0.35cm,center]{frametitle}
    {\usebeamerfont{frametitle}\insertframetitle\par}
    \ifx\insertframesubtitle\@empty\else
      \vspace{0.15cm}
      {\usebeamerfont{framesubtitle}\usebeamercolor[fg]{framesubtitle}\insertframesubtitle\par}
    \fi
  \end{beamercolorbox}
}
\makeatother

% ---------- Margins ----------
\setbeamersize{
  text margin left=2em,
  text margin right=2em
}

% ---------- Section outline frames ----------
\AtBeginSection[]{
  \begin{frame}{Outline}
    \tableofcontents[currentsection,
                     currentsubsection,
                     hideothersubsections]
  \end{frame}
}

% =============================================================================
% Colors — Okabe-Ito colorblind-friendly palette
% =============================================================================

\definecolor{black}{HTML}{000000}
\definecolor{orange}{HTML}{E69F00}
\definecolor{skyblue}{HTML}{56B4E9}
\definecolor{green}{HTML}{009E73}
\definecolor{yellow}{HTML}{F0E442}
\definecolor{blue}{HTML}{0072B2}
\definecolor{vermillion}{HTML}{D55E00}
\definecolor{purple}{HTML}{CC79A7}

% Semantic aliases (mapped to Okabe-Ito)
\definecolor{softbg}{HTML}{F5F5F5}
\definecolor{lightgray}{gray}{0.65}

% =============================================================================
% Box Environments
% =============================================================================

% --- Key points (orange background) ---
\newtcolorbox{keybox}{
  colback=orange!12,
  colframe=orange,
  fonttitle=\bfseries,
  boxrule=0.5pt,
  arc=2pt,
  left=6pt,
  right=6pt,
  top=4pt,
  bottom=4pt
}

% --- Highlights (orange left accent) ---
\newtcolorbox{highlightbox}{
  colback=white,
  colframe=orange,
  fonttitle=\bfseries,
  boxrule=0pt,
  leftrule=3pt,
  arc=0pt,
  left=6pt,
  right=6pt,
  top=4pt,
  bottom=4pt
}

% --- Definitions (blue border, titled) ---
\newtcolorbox{definitionbox}[1][]{
  colback=blue!5,
  colframe=blue,
  fonttitle=\bfseries,
  title={#1},
  boxrule=0.5pt,
  arc=2pt,
  left=6pt,
  right=6pt,
  top=4pt,
  bottom=4pt
}

% --- Methods (sky blue background) ---
\newtcolorbox{methodbox}{
  colback=skyblue!8,
  colframe=skyblue!70,
  fonttitle=\bfseries,
  boxrule=0.5pt,
  arc=2pt,
  left=6pt,
  right=6pt,
  top=4pt,
  bottom=4pt
}

% --- Results (green border) ---
\newtcolorbox{resultbox}{
  colback=green!5,
  colframe=green,
  fonttitle=\bfseries,
  boxrule=0.5pt,
  arc=2pt,
  left=6pt,
  right=6pt,
  top=4pt,
  bottom=4pt
}

% --- Quotes (purple left rule) ---
\newtcolorbox{quotebox}{
  colback=softbg,
  colframe=purple,
  fonttitle=\itshape,
  boxrule=0pt,
  leftrule=2pt,
  arc=0pt,
  left=6pt,
  right=6pt,
  top=4pt,
  bottom=4pt
}

% =============================================================================
% Math Commands
% =============================================================================

\newcommand{\E}{\mathrm{E}}
\newcommand{\Var}{\mathrm{Var}}
\newcommand{\Cov}{\mathrm{Cov}}
\newcommand{\R}{\mathbb{R}}
\newcommand{\suminfty}[1]{\sum_{#1=0}^\infty}
\newcommand{\Lcal}{\mathcal{L}}
\newcommand{\prtl}[2]{\frac{\partial #1}{\partial #2}}
\renewcommand{\div}[2]{\frac{d #1}{d #2}}
\newcommand{\suchthat}{\text{ s.t. }}
\newcommand{\wrt}{\text{ w.r.t. }}
\newcommand{\red}[1]{\textcolor{vermillion}{#1}}
\newcommand{\blue}[1]{\textcolor{blue}{#1}}

% =============================================================================
% Economics Shortcuts
% =============================================================================

\newcommand{\OC}{\text{Opportunity Cost}}
\newcommand{\DWL}{\text{DWL}}
\newcommand{\pmark}{\textcolor{resultgreen}{\checkmark}}
\newcommand{\xmark}{\textcolor{darkred}{\texttimes}}

% Compile: cd output/Slides && TEXINPUTS=../../codes/Preambles:$TEXINPUTS xelatex file.tex
% =============================================================================

% ---------- Packages ----------
% Fonts
\usepackage{lmodern}
\usepackage[utf8]{inputenc}
\usepackage[normalem]{ulem}

% Math
\usepackage{amsmath,amssymb,amsthm}

% Figures and tables
\usepackage{graphicx}
\usepackage{array,tabularx,booktabs,threeparttable}

% TikZ and pgfplots
\usepackage{pgfplots}
\pgfplotsset{compat=1.18}
\usepackage{tikz}
\usetikzlibrary{positioning,arrows.meta,calc}

% Colored boxes
\usepackage[most]{tcolorbox}

% Miscellaneous
\usepackage{transparent}
\usepackage{setspace}
\usepackage[nodayofweek]{datetime}
\usepackage{hyperref}

% ---------- Theme ----------
\usetheme{CambridgeUS}
\usecolortheme{dolphin}
\usepackage{appendixnumberbeamer}

% ---------- Itemize ----------
\setbeamercovered{invisible}
\setbeamertemplate{itemize items}[default]
\setbeamertemplate{enumerate items}[default]

\setbeamertemplate{itemize item}[triangle]
\setbeamertemplate{itemize subitem}[triangle]
\setbeamertemplate{itemize subsubitem}{-}

% ---------- Headline / Footline ----------
\setbeamertemplate{headline}{}

\setbeamertemplate{footline}{
  \leavevmode
  \hbox{
    \begin{beamercolorbox}[wd=\paperwidth,ht=2.5ex,dp=1.25ex,leftskip=1em,rightskip=1em]{footline}
      \hfill \insertframenumber
    \end{beamercolorbox}
  }
}

\setbeamertemplate{navigation symbols}{}

% ---------- Frame title (centered) ----------
\makeatletter
\setbeamertemplate{frametitle}{
  \nointerlineskip
  \begin{beamercolorbox}[wd=\paperwidth,sep=0.35cm,center]{frametitle}
    {\usebeamerfont{frametitle}\insertframetitle\par}
    \ifx\insertframesubtitle\@empty\else
      \vspace{0.15cm}
      {\usebeamerfont{framesubtitle}\usebeamercolor[fg]{framesubtitle}\insertframesubtitle\par}
    \fi
  \end{beamercolorbox}
}
\makeatother

% ---------- Margins ----------
\setbeamersize{
  text margin left=2em,
  text margin right=2em
}

% ---------- Section outline frames ----------
\AtBeginSection[]{
  \begin{frame}{Outline}
    \tableofcontents[currentsection,
                     currentsubsection,
                     hideothersubsections]
  \end{frame}
}

% =============================================================================
% Colors — Okabe-Ito colorblind-friendly palette
% =============================================================================

\definecolor{black}{HTML}{000000}
\definecolor{orange}{HTML}{E69F00}
\definecolor{skyblue}{HTML}{56B4E9}
\definecolor{green}{HTML}{009E73}
\definecolor{yellow}{HTML}{F0E442}
\definecolor{blue}{HTML}{0072B2}
\definecolor{vermillion}{HTML}{D55E00}
\definecolor{purple}{HTML}{CC79A7}

% Semantic aliases (mapped to Okabe-Ito)
\definecolor{softbg}{HTML}{F5F5F5}
\definecolor{lightgray}{gray}{0.65}

% =============================================================================
% Box Environments
% =============================================================================

% --- Key points (orange background) ---
\newtcolorbox{keybox}{
  colback=orange!12,
  colframe=orange,
  fonttitle=\bfseries,
  boxrule=0.5pt,
  arc=2pt,
  left=6pt,
  right=6pt,
  top=4pt,
  bottom=4pt
}

% --- Highlights (orange left accent) ---
\newtcolorbox{highlightbox}{
  colback=white,
  colframe=orange,
  fonttitle=\bfseries,
  boxrule=0pt,
  leftrule=3pt,
  arc=0pt,
  left=6pt,
  right=6pt,
  top=4pt,
  bottom=4pt
}

% --- Definitions (blue border, titled) ---
\newtcolorbox{definitionbox}[1][]{
  colback=blue!5,
  colframe=blue,
  fonttitle=\bfseries,
  title={#1},
  boxrule=0.5pt,
  arc=2pt,
  left=6pt,
  right=6pt,
  top=4pt,
  bottom=4pt
}

% --- Methods (sky blue background) ---
\newtcolorbox{methodbox}{
  colback=skyblue!8,
  colframe=skyblue!70,
  fonttitle=\bfseries,
  boxrule=0.5pt,
  arc=2pt,
  left=6pt,
  right=6pt,
  top=4pt,
  bottom=4pt
}

% --- Results (green border) ---
\newtcolorbox{resultbox}{
  colback=green!5,
  colframe=green,
  fonttitle=\bfseries,
  boxrule=0.5pt,
  arc=2pt,
  left=6pt,
  right=6pt,
  top=4pt,
  bottom=4pt
}

% --- Quotes (purple left rule) ---
\newtcolorbox{quotebox}{
  colback=softbg,
  colframe=purple,
  fonttitle=\itshape,
  boxrule=0pt,
  leftrule=2pt,
  arc=0pt,
  left=6pt,
  right=6pt,
  top=4pt,
  bottom=4pt
}

% =============================================================================
% Math Commands
% =============================================================================

\newcommand{\E}{\mathrm{E}}
\newcommand{\Var}{\mathrm{Var}}
\newcommand{\Cov}{\mathrm{Cov}}
\newcommand{\R}{\mathbb{R}}
\newcommand{\suminfty}[1]{\sum_{#1=0}^\infty}
\newcommand{\Lcal}{\mathcal{L}}
\newcommand{\prtl}[2]{\frac{\partial #1}{\partial #2}}
\renewcommand{\div}[2]{\frac{d #1}{d #2}}
\newcommand{\suchthat}{\text{ s.t. }}
\newcommand{\wrt}{\text{ w.r.t. }}
\newcommand{\red}[1]{\textcolor{vermillion}{#1}}
\newcommand{\blue}[1]{\textcolor{blue}{#1}}

% =============================================================================
% Economics Shortcuts
% =============================================================================

\newcommand{\OC}{\text{Opportunity Cost}}
\newcommand{\DWL}{\text{DWL}}
\newcommand{\pmark}{\textcolor{resultgreen}{\checkmark}}
\newcommand{\xmark}{\textcolor{darkred}{\texttimes}}


% ---------- Title ----------
\title{\textbf{Economic Decision-Making and Everyday Choices}}
\subtitle{Pathways to Arts \& Humanities Summer Scholars Program}
\author{Instructor: Tudor Schlanger}
\date{Summer 2026}

\begin{document}

% =============================================================================
% Title
% =============================================================================
\begin{frame}
  \titlepage
\end{frame}

% =============================================================================
\section{What Is Economics?}
% =============================================================================

\begin{frame}{What Do Economists Study?}

  Economists recognize a fundamental problem: \textbf{Scarcity}. \pause 

  \begin{definitionbox}[Economics]
    Economics is the study of how individuals and societies allocate scarce resources. 
  \end{definitionbox} \pause 

  What does it mean to allocate resources? Do resources just fall from a tree?  \pause 
  
  \begin{definitionbox}[Economics]
    Economics is the study of how individuals and societies \textbf{produce, exchange, distribute and consume} scarce resources, such as \textbf{goods and services}.  
  \end{definitionbox}


\end{frame}

\begin{frame}{What Do Economists Study?}

  \begin{definitionbox}[Economics]
    Economics is the study of how individuals and societies \textbf{produce, exchange, distribute and consume} scarce resources, such as \textbf{goods and services}.  
  \end{definitionbox}

  Economics has grown to study many areas of the economy: 

  \vspace{0.3em}
  \begin{columns}[T]
    \begin{column}{0.3\textwidth}
      \begin{itemize}
        \item Environment \& climate
        \item Health care
        \item Education
      \end{itemize}
    \end{column}
    \begin{column}{0.3\textwidth}
      \begin{itemize}
        \item Trade \& tariffs
        \item Finance \& markets
        \item Labor \& employment
      \end{itemize}
    \end{column}
    \begin{column}{0.3\textwidth}
      \begin{itemize}
        \item Industrial organization
        \item Development
        \item Public policy
      \end{itemize}
    \end{column}
  \end{columns} \pause 

  \vspace{2em}
  Do all these fall under the definition above? 

\end{frame}

\begin{frame}{What Do Economists Study?}

  \begin{definitionbox}[Economics]
    Economics is the study of how individuals and societies \textbf{produce, exchange, distribute and consume} scarce resources, such as \textbf{goods and services}.  
  \end{definitionbox}

  Economics has grown to study many areas of the economy: 

  \vspace{0.3em}
  \begin{columns}[T]
    \begin{column}{0.3\textwidth}
      \begin{itemize}
        \item Environment \& climate
        \item Health care
        \item \textbf{Education}
      \end{itemize}
    \end{column}
    \begin{column}{0.3\textwidth}
      \begin{itemize}
        \item Trade \& tariffs
        \item Finance \& markets
        \item Labor \& employment
      \end{itemize}
    \end{column}
    \begin{column}{0.3\textwidth}
      \begin{itemize}
        \item Industrial organization
        \item Development
        \item Public policy
      \end{itemize}
    \end{column}
  \end{columns}  \pause 

    \vspace{2em}

  Gary Becker, Nobel Prize winner,  argued that an individual’s investment in education is similar to a company’s investment in new machinery or equipment.  \pause 

    \vspace{2em}

  Time is a scarce resource too! 


\end{frame}

\begin{frame}{Scarcity creates trade-offs}


  A key part of economics is about studying \textbf{trade-offs}:
  \begin{itemize}
    \item What do you \textcolor{green}{gain} with each decision?
    \item What do you \textcolor{vermillion}{give up}?
  \end{itemize}

  \vspace{0.8em}
  Every decision---from what to eat for lunch to how a government spends
  billions---comes with \textbf{costs} and \textbf{benefits}.

  We will consider two costs that economists use all the time:
  \begin{enumerate}
    \item Opportunity Cost
    \item Sunk Cost
  \end{enumerate}

\end{frame}

% =============================================================================
\section{Opportunity Cost}
% =============================================================================

% ---------- Definition ----------
\begin{frame}{Opportunity Cost}

  \begin{definitionbox}[Opportunity Cost]
    The \textit{net} value of the \textit{next best alternative} when you make a decision.
  \end{definitionbox}

  \vspace{0.8em}
  \begin{keybox}
    \textbf{You have a free evening.}\\
    Option A: study for a test.\\
    Option B: go to a movie (for free).\\
    You choose Option A (study). What is the opportunity cost?
  \end{keybox}

\end{frame}

% ---------- Opening Example: Question ----------
\begin{frame}{Opportunity Cost}

  \begin{definitionbox}[Opportunity Cost]
    The \textit{net} value of the \textit{next best alternative} when you make a decision.
  \end{definitionbox}

  \vspace{0.8em}
  \begin{keybox}
    \textbf{You have a free evening.}\\
    Option A: study for a test.\\
    Option B: go to a movie (for free).\\
    Option C: go to the park.\\[0.3em]
    You choose Option A (study). What is the opportunity cost?
  \end{keybox}

  \vspace{1.5em}
  \begin{center}
    \Large\textcolor{blue}{Think for 30 seconds\ldots}
  \end{center}

\end{frame}

% ---------- Opening Example: Answer ----------
\begin{frame}{Opportunity Cost}

  \begin{definitionbox}[Opportunity Cost]
    The \textit{net} value of the \textit{next best alternative} when you make a decision.
  \end{definitionbox}

  \vspace{0.8em}
  \begin{keybox}
    \textbf{You have a free evening.}\\
    Option A: study for a test.\\
    Option B: go to a movie (for free).\\
    Option C: go to the park.\\[0.3em]
    You choose Option A (study). What is the opportunity cost?
  \end{keybox}

  \vspace{0.5em}
  \begin{resultbox}
    Let's assume you like the movies more. Then, the opportunity cost of studying is the \textbf{value of going to the
    movies} (Option B)---the next best alternative you gave up, not all
    alternatives combined.
  \end{resultbox}

\end{frame}


\begin{frame}{``Free'' Things}

  \begin{keybox}
    A friend invites Maya to a \textbf{free} outdoor concert---three
    hours, no admission fee.\\ 
    
    Instead, she could work her shift at the caf\'{e} (\$12/hr $\times$ 3 hrs
    = \$36).\\ Assume there are no costs to her working there (e.g. bus tickets). What is her opportunity cost? 
  \end{keybox}

  \vspace{1.5em}
  \begin{center}
    \Large\textcolor{blue}{Discuss with your group for 30 seconds \ldots}
  \end{center}

\end{frame}

\begin{frame}{``Free'' Things}

  \begin{keybox}
    A friend invites Maya to a \textbf{free} outdoor concert---three
    hours, no admission fee.\\ 
    
    Instead, she could work her shift at the caf\'{e} (\$12/hr $\times$ 3 hrs
    = \$36).\\ Assume there are no costs to her working there (e.g. bus tickets). What is her opportunity cost? 
  \end{keybox}

  \vspace{0.5em}
  \begin{resultbox}
    The ``free'' concert costs Maya \textbf{\$36} in foregone earnings. 
  \end{resultbox}

  \vspace{0.5em}
  \begin{highlightbox}
    ``Free'' is one of the most misleading words in everyday life.
    When something costs no money, it almost always costs \textbf{time}---and
    time has value.
  \end{highlightbox}

\end{frame}

\begin{frame}{Direct vs Opportunity Costs}

  \begin{keybox}
    Maya goes to a movie. Ticket: \$14.\ \ Snacks: \$9.\
    Three hours of her time, during which she could have earned \$36.\\
    What is the direct cost of the movie? What is the opportunity cost?
  \end{keybox}

  \vspace{1.5em}
  \begin{center}
    \Large\textcolor{blue}{Think for 30 seconds\ldots}
  \end{center}

\end{frame}

\begin{frame}{Direct vs Opportunity Costs}

  \begin{keybox}
    Maya goes to a movie. Ticket: \$14.\ \ Snacks: \$9.\
    Three hours---she could have earned \$36.
  \end{keybox}

  \vspace{0.5em}
  \begin{resultbox}
    Direct cost = \$14 + \$9  = \$23 \\
    Opportunity cost = \$14 + \$9 + \$36 = \textbf{\$59}.
  \end{resultbox}

  \vspace{0.5em}
  \begin{highlightbox}
  Real decisions combine \textbf{direct costs} (what you pay out of
  pocket) and the value of what you give up. The total is the \textbf{opportunity cost}. 
  \end{highlightbox}

\end{frame}

% ---------- Example 6: College vs. Work [Groups] ----------
\begin{frame}{Real world example: College vs.\ Full-Time Work}

  \begin{keybox}
    After high school, Maya considers college vs.\ starting full-time
    work.\\ Tuition: \textasciitilde\$20{,}000/year.\ \ Foregone wages
    from a full-time job: \textasciitilde\$35{,}000/year.\\
    What is the opportunity cost of one year of college?
  \end{keybox}

\end{frame}

% % ---------- Transition ----------
% \begin{frame}{Everyone Has a Different Opportunity Cost}

%   Suppose you're given \$100. What is the opportunity cost of
%   \textbf{saving} it?

%   \vspace{0.5em}
%   \begin{columns}[c]
%     \begin{column}{0.3\textwidth}
%       \centering
%       \textcolor{blue}{\large Person A}\\[0.3em]
%       Needs groceries badly.\\
%       Opportunity cost:\\
%       \emph{going hungry}
%     \end{column}
%     \begin{column}{0.3\textwidth}
%       \centering
%       \textcolor{orange}{\large Person B}\\[0.3em]
%       Wants concert tickets.\\
%       Opportunity cost:\\
%       \emph{missing a great show}
%     \end{column}
%     \begin{column}{0.3\textwidth}
%       \centering
%       \textcolor{green}{\large Person C}\\[0.3em]
%       Already has savings.\\
%       Opportunity cost:\\
%       \emph{very little}
%     \end{column}
%   \end{columns}

%   \vspace{0.8em}
%   \begin{highlightbox}
%     People face different trade-offs because they have different
%     \textbf{preferences}, \textbf{resources}, and \textbf{alternatives}.
%     That's why different people make different choices.
%   \end{highlightbox}

% \end{frame}

\begin{frame}{Making Good Choices}

  \begin{resultbox}
    A ``good'' economic choice is one where the \textbf{benefits}
    are greater than the \textbf{opportunity cost}.
  \end{resultbox}

  \vspace{0.6em}
  \begin{itemize}
    \item List what you gain.
    \item List what you give up.
    \item Compare---and choose.
  \end{itemize}

  \vspace{0.8em}
  \textbf{Let's practice this for the college example!}


\end{frame}

% ---------- Example 6: College vs. Work [Groups] ----------
\begin{frame}{Example 6: College vs.\ Full-Time Work}

  After high school, Maya considers college vs.\ starting full-time work.

  \vspace{0.5em}
  \centering
  \begin{tabular}{l c c}
    \toprule
    & \textbf{College} & \textbf{No College (work full-time)} \\
    \midrule
    Value of college experience    & \$10,000     & --- \\
    Earnings (first 5 yrs)         & \$0          & \$35,000/yr \\
    Earnings (next 35 yrs)         & \$100,000/yr & \$70,000/yr \\
    Tuition  (first 5 yrs)         & \$20,000/yr  & --- \\
    \bottomrule
  \end{tabular}

  \vspace{1em}
  \begin{center}
    \Large\textcolor{blue}{Discuss with your group for 2 minutes\ldots}
  \end{center}

\end{frame}

\begin{frame}{Example 6: College vs.\ Full-Time Work}

  After high school, Maya considers college vs.\ starting full-time work.

  \vspace{0.5em}
  \centering
  \begin{tabular}{l c c}
    \toprule
    & \textbf{College} & \textbf{No College (work full-time)} \\
    \midrule
    Value of college experience    & \$10,000     & --- \\
    Earnings (first 5 yrs)         & \$0          & \$35,000/yr \\
    Earnings (next 35 yrs)         & \$100,000/yr & \$70,000/yr \\
    Tuition  (first 5 yrs)         & \$20,000/yr  & --- \\
    \bottomrule
  \end{tabular}

  \vspace{0.5em}
  \begin{resultbox}
    \textbf{College:}\\
    Benefits = \$10{,}000 + (35 $\times$ \$100{,}000) = \$3{,}510{,}000\\
    Opportunity Costs = (5 $\times$ \$20{,}000) +  (5 $\times$ \$35{,}000) + (35 $\times$ \$70{,}000) = \$2{,}725{,}000 \\
    Net Value = Benefits - Opportunity Costs = \textcolor{green}{\$785{,}000} > 0 
    \vspace{0.3em}
  \end{resultbox}

  \begin{keybox}
    The only thing we need to make sure is that the net value is positive! 
  \end{keybox}


\end{frame}



% =============================================================================
\section{Sunk Costs}
% =============================================================================

\begin{frame}{What Is a Sunk Cost?}

  \begin{definitionbox}[Sunk Cost]
    A cost you \textbf{cannot recover}---money, time, or effort
    already spent and gone forever.
  \end{definitionbox}

  \vspace{0.8em}
  \textbf{Everyday examples:}
  \begin{itemize}
    \item You've waited 20 minutes in a terrible line. Do you stay or leave?
    \item You paid \$15 for a movie that's awful. Do you walk out?
    \item You bought an expensive gym membership you never use.
  \end{itemize}

\end{frame}

\begin{frame}{The Sunk Cost Fallacy}

  \begin{definitionbox}[Sunk Cost Fallacy]
    Letting past, unrecoverable costs influence your \textbf{future}
    decisions---staying on a path just because you've already invested in it,
    even when it no longer makes sense.
  \end{definitionbox}

  \vspace{0.5em}
  \textbf{Classic sign:}
  \begin{quotebox}
    ``I can't stop now---I've already put so much time/money into this!''
  \end{quotebox}

  \vspace{0.3em}
  The time and money are \textbf{already gone} whether you stay or leave.
  The only question that matters is: \emph{What should I do from here?}

\end{frame}

\begin{frame}{Alex's Bitcoin Story}

  \begin{itemize}
    \item \textbf{November 2021:} Alex buys \$10{,}000 of Bitcoin at
          \$68{,}789 per coin.
    \vspace{0.4em}
    \item \textbf{November 2022:} Bitcoin crashes to \$15{,}800.\\
          Alex's \$10{,}000 is now worth \textbf{\textasciitilde\$2{,}300}.
          That's a \textcolor{vermillion}{\$7{,}700 loss}.
    \vspace{0.4em}
    \item Alex thinks: \emph{``I can't sell now---I'd lock in the loss!
          If I hold on, maybe it'll go back up.''}
  \end{itemize}

  \vspace{0.5em}
  \begin{highlightbox}
    The \$7{,}700 loss \textbf{already happened}. Whether Alex sells or
    holds, that money is gone. The real question is about the future.
  \end{highlightbox}

\end{frame}

\begin{frame}{The Right Question to Ask}

  \begin{methodbox}
    \textbf{The Reframe:} ``If I had \$2{,}300 in cash right now---would I
    buy Bitcoin with it?''
  \end{methodbox}

  \vspace{0.5em}
  \begin{columns}[T]
    \begin{column}{0.45\textwidth}
      \centering
      \textcolor{green}{\large If YES $\rightarrow$ Hold}\\[0.3em]
      You believe Bitcoin will grow from here.
      Holding is based on the \textbf{future}, not the past.
    \end{column}
    \begin{column}{0.45\textwidth}
      \centering
      \textcolor{vermillion}{\large If NO $\rightarrow$ Sell}\\[0.3em]
      You're only holding because of the past.
      That's the sunk cost fallacy.
    \end{column}
  \end{columns}

\end{frame}

\begin{frame}{How to Avoid the Sunk Cost Fallacy}

  \begin{enumerate}
    \item \textbf{Ignore the sunk costs}---what you've already spent
          and cannot recover.
    \vspace{0.4em}
    \item \textbf{Look at what you have now}---what are your current
          options \emph{today}?
    \vspace{0.4em}
    \item \textbf{Decide fresh}---base your choice only on future costs
          and benefits.
  \end{enumerate}

  \vspace{0.6em}
  This applies to everything: bad movies, classes you hate, relationships,
  business projects, gym memberships.

  \vspace{0.6em}
  \begin{quotebox}
    ``If you find yourself in a hole, the first thing to do is
    \textbf{stop digging}.''
  \end{quotebox}

\end{frame}

% =============================================================================
\section{Policy Debate: Cell Phone Ban}
% =============================================================================

\begin{frame}{Policy Debate}

  \begin{keybox}
    \textbf{Economic analysis:} Is a ``phone-free school day'' a good policy?
  \end{keybox}

  \vspace{0.8em}
  \textbf{Question:} Should Connecticut pass a \textbf{bell-to-bell} phone
  ban for all public schools?

  \vspace{0.5em}
  Three stakeholder groups will analyze the question:
  \begin{itemize}
    \item \textcolor{blue}{Students}
    \item \textcolor{orange}{Teachers}
    \item \textcolor{green}{Government officials}
  \end{itemize}

  \vspace{0.5em}
  Each group will identify \textbf{benefits}, \textbf{costs},
  \textbf{opportunity costs}, and any \textbf{sunk costs}---then take
  a policy stance.

\end{frame}

\begin{frame}{Your Task}

  For your assigned group, answer:
  \vspace{0.3em}

  \begin{enumerate}
    \item What is your \textbf{goal}? What are your priorities and
          preferences?
    \vspace{0.3em}
    \item What are the \textcolor{green}{benefits} of a ban?
          \begin{itemize}
            \item List at least one \textbf{opportunity cost} of
                  \emph{not} banning phones.
            \item Who gets the benefits?
          \end{itemize}
    \vspace{0.3em}
    \item What are the \textcolor{vermillion}{costs} of a ban?
          \begin{itemize}
            \item List at least one \textbf{opportunity cost} of
                  banning phones.
            \item List at least one \textbf{sunk cost} in this situation.
          \end{itemize}
    \vspace{0.3em}
    \item What is your \textbf{policy stance}? Defend it.
  \end{enumerate}

\end{frame}

% =============================================================================
\section{Recap}
% =============================================================================

\begin{frame}{Today's Key Takeaways}

  \begin{definitionbox}[Opportunity Cost]
    The value of the \textbf{next best alternative} when you make a choice.
  \end{definitionbox}

  \vspace{0.3em}

  \begin{definitionbox}[Sunk Cost]
    A cost you \textbf{cannot recover}---money, time, or effort already spent.
  \end{definitionbox}

  \vspace{0.3em}

  \begin{definitionbox}[Sunk Cost Fallacy]
    Considering sunk costs when deciding about the \textbf{future}---staying on
    a path only because of what you've already invested.
  \end{definitionbox}

  \vspace{0.4em}
  \begin{highlightbox}
    \textbf{Good economic thinking:} Compare benefits and opportunity costs.
    Ignore sunk costs. Decide based on what matters \emph{going forward}.
  \end{highlightbox}

\end{frame}

\end{document}
